% !TeX program = xelatex
\documentclass[9pt,letterpaper]{extarticle}

\usepackage[margin=0.55in]{geometry}
\usepackage{fontspec}
\usepackage{xcolor}
\usepackage{hyperref}
\usepackage{enumitem}
\usepackage{tabularx}
\usepackage{array}
\usepackage{titlesec}
\usepackage{microtype}
\usepackage{ragged2e}
\usepackage{needspace}
\usepackage{lastpage}
\usepackage{fancyhdr}

\setmainfont{Noto Serif}
\setsansfont{Noto Sans}
\newfontfamily\chinesefont{Noto Serif CJK SC}

\definecolor{accent}{HTML}{243B63}
\definecolor{textgray}{HTML}{333333}
\definecolor{muted}{HTML}{666666}
\definecolor{rulegray}{HTML}{D6DCE8}

\hypersetup{
  colorlinks=true,
  urlcolor=accent,
  linkcolor=accent
}

\pagestyle{fancy}
\fancyhf{}
\renewcommand{\headrulewidth}{0pt}
\cfoot{\sffamily\footnotesize\color{muted}Xujiang Tang - Curriculum Vitae - Page \thepage\ of \pageref*{LastPage}}

\setlength{\parindent}{0pt}
\setlength{\parskip}{0pt}
\renewcommand{\arraystretch}{1.08}
\setlist[itemize]{leftmargin=1.05em, topsep=1pt, itemsep=0.6pt, parsep=0pt}
\setlist[enumerate]{leftmargin=1.25em, topsep=1pt, itemsep=1.1pt, parsep=0pt}
\linespread{0.98}

\titleformat{\section}
  {\Needspace{5\baselineskip}\sffamily\large\bfseries\color{accent}}
  {}{0pt}{}[\vspace{-0.55em}\color{rulegray}\titlerule]
\titlespacing*{\section}{0pt}{0.62em}{0.35em}

\newcommand{\name}[1]{\begin{center}{\fontsize{23}{26}\selectfont\bfseries\color{accent}#1}\end{center}}
\newcommand{\tagline}[1]{\begin{center}{\sffamily\small\color{textgray}#1}\end{center}\vspace{-0.25em}}
\newcommand{\contact}[1]{\begin{center}{\small\color{muted}#1}\end{center}}

\newcolumntype{L}{>{\RaggedLeft\arraybackslash\color{muted}\sffamily\footnotesize}p{0.18\linewidth}}
\newcolumntype{R}{>{\RaggedRight\arraybackslash}p{0.77\linewidth}}
\newcommand{\entry}[4]{%
  \begin{tabularx}{\linewidth}{@{}L@{\hspace{0.035\linewidth}}R@{}}
  #1 & {\bfseries #2}\hfill {\sffamily\footnotesize\color{muted}#3}\\[-0.15em]
     & #4
  \end{tabularx}\vspace{0.18em}
}
\newcommand{\paper}[2]{\item #1. ``#2.''}
\newcommand{\skill}[2]{\textbf{#1:} #2\par\vspace{0.05em}}

\begin{document}

\name{Xujiang Tang {\chinesefont 唐煦江}}
\tagline{Theoretical Mathematics Researcher \; | \; Machine Learning Theory Enthusiast}
\tagline{Undergraduate Student in Mathematics and Applied Mathematics, Yangtze University \; | \; Visiting Student, Great Bay University}
\contact{\href{mailto:txj_262538@163.com}{txj\_262538@163.com} \;\textbar\; \href{https://txj2006.github.io}{txj2006.github.io} \;\textbar\; \href{https://github.com/TXJ2006}{GitHub} \;\textbar\; \href{https://orcid.org/0009-0008-1127-6420}{ORCID} \;\textbar\; \href{https://www.linkedin.com/in/\%E7\%85\%A6\%E6\%B1\%9F-\%E5\%94\%90-7a1513392}{LinkedIn}}
\vspace{0.45em}
\color{rulegray}\hrule height 0.9pt\color{textgray}
\vspace{0.55em}

\section{Research Interests}
\begin{itemize}
  \item \textbf{Algebraic topology and equivariant homotopy theory:} fixed point methods, geometric fixed point tomography, Bökstedt periodicity, and stable homotopy-theoretic structures.
  \item \textbf{Finite group theory:} Quillen's p-subgroup conjecture, subgroup complexes, centralizer structures, and FRUL-WOS sphere methods.
  \item \textbf{Formal proof:} Lean, Mathlib, and the formalization of advanced mathematics in algebra and topology.
  \item \textbf{Machine learning theory:} generalization, learnability, out-of-distribution prediction, and Koopman/operator-theoretic methods for long-sequence data.
\end{itemize}

\section{Education}
\entry{Expected Jun. 2027}{Yangtze University College of Arts and Sciences}{Jingzhou, China}{B.Sc. in Mathematics and Applied Mathematics. Coursework: mathematical analysis, linear algebra, probability and statistics, differential equations, computational modeling, optimization, numerical computation, data structures, Python, R, MATLAB. Marks: Optimization Theory 97/100; Numerical Computation Methods 96/100.}

\section{Visiting Position and Training}
\entry{2025--present}{Visiting Student, Great Bay University}{}{Research in theoretical mathematics and machine learning theory.}
\entry{2026}{Zhejiang University Lean Formalization Summer School}{}{Formal proof, Lean, Mathlib, and computer-assisted theorem proving.}
\entry{Winter training}{Neural Networks and Python Implementation}{Grade: A}{Short winter training with Prof. Mark Vogelsberger. Completed the course and received an A grade.}

\section{Research Experience}
\entry{2025--present}{Finite Group Theory and Quillen's p-Subgroup Conjecture}{}{Centralizer energy, FRUL-WOS spheres, and subgroup-complex methods for Quillen's p-subgroup conjecture.}
\entry{2025--present}{Algebraic Topology and Equivariant Homotopy Theory}{}{Geometric fixed point tomography and spoke Bökstedt periodicity.}
\entry{2026--present}{Formal Proof and Lean}{}{Lean formalization, proof engineering, and formalization of advanced mathematical arguments.}
\entry{2024--present}{Theoretical Machine Learning and Optimization}{}{Learnability, generalization guarantees, out-of-distribution prediction, and optimization-theoretic aspects of learning algorithms.}
\entry{2023--present}{Stochastic Processes, Dynamical Systems, and Quantitative Finance}{}{Stochastic processes, rough volatility, MCMC simulation, Koopman/operator-theoretic dynamics, and robust financial decision problems.}
\entry{2023--2025}{Epidemic Prevention and Population Mobility Modeling, Jingzhou}{}{Mathematical modeling and statistical surveys on epidemic prevention and population mobility in Jingzhou.}
\entry{2025--2026 H1}{Spatial Omics Data Processing}{}{Spatial omics data processing with Dayu Hu, Northeastern University.}

\newpage
\section*{Research Experience (continued)}
\entry{2025--2026 H2}{Epilepsy and Long-Sequence Data Research}{}{Epilepsy-related long-sequence modeling in Prof. Huaguang Gu's group, focusing on Koopman operator methods.}

\section{Publications and Manuscripts}
\textbf{Published and Accepted Papers}
\begin{enumerate}
  \paper{Q. Li and \textbf{X. Tang}}{Robust Optimal Reinsurance and Investment with Inflation Risk: A Game-Theoretic Approach and Explicit Solutions} \emph{AIMS Mathematics}, 11(3), 7330--7352, 2026. Published. DOI: \href{https://doi.org/10.3934/math.2026302}{10.3934/math.2026302}.
  \paper{\textbf{X. Tang et al.}}{Structure-preserving Koopman Predictive Control for Memristive Neural Dynamics: Input-exact and Commutator-defect Lifting} \emph{Biological Cybernetics}. Accepted. \href{https://doi.org/10.21203/rs.3.rs-9200134/v1}{Public preprint}.
\end{enumerate}

\textbf{Under Review, Submitted, and Manuscripts}
\begin{enumerate}
  \paper{\textbf{X. Tang et al.}}{The Curvature Filter in Bilevel Optimization: Implicit Differentiation Without Nondegeneracy} \emph{Journal of Machine Learning Research}, Theory and Methods. Submitted.
  \paper{\textbf{X. Tang}, R. Guan, and Z. Wang}{Moment Order in Unsupervised Direction Learning} \emph{Journal of Machine Learning Research}. Submitted.
  \paper{Y. Fan, \textbf{X. Tang}, and R. Guan}{Deep Hedging under Rough Volatility: A Fractional Kernel Embedding Approach with Optimal Convergence Rate} \emph{Journal of Computational and Applied Mathematics}. Under review.
  \paper{\textbf{X. Tang et al.}}{Spectral Network Determinants of Seizure-like Synchronization and Spread in Coupled Hindmarsh-Rose Brain Models} \emph{Cognitive Neurodynamics}. Under review.
  \paper{\textbf{X. Tang et al.}}{Trajectory-Level Out-of-Distribution Success Bounds for Deep Autoregressive Models} Manuscript, 2026.
\end{enumerate}

\section{Collaborators and Mentors}
\begin{itemize}
  \item \href{https://engineering.jhu.edu/faculty/william-redman/}{William Redman}, Johns Hopkins University.
  \item \href{https://qiaoziyue.com/}{Ziyue Qiao}, Great Bay University.
\end{itemize}

\section{Software Copyrights}
\begin{itemize}
  \item Software copyright in financial systems.
  \item Software copyright in unmanned swarm systems.
\end{itemize}

\section{Honors, Competitions, and Activities}
\begin{itemize}
  \item Kaggle Silver Medal, Jigsaw Toxic Comment Classification Challenge.
  \item Regional Gold Medal, WorldQuant International Quant Championship.
  \item Third Prize, UCAS Graduate AI Forum.
  \item Mathematical Modeling Competition, Hua Zhong Cup.
  \item National College Students Statistical Modeling Competition.
  \item Multiple national second prizes in mathematical modeling competitions.
  \item Reviewer, ICLR 2026 Workshop.
\end{itemize}

\section{Professional Experience}
\entry{Jun.--Aug. 2024}{Peking University - Open-Source LLM Group}{}{Contributed to LLM deployment and retrieval-augmented generation systems.}
\entry{Jul.--Sep. 2024}{Deloitte}{}{Professional experience in consulting workflows, teamwork, and communication.}
\entry{Jul.--Sep. 2024}{Guolian Securities}{}{Applied statistical methods to securities research and quantitative workflows.}

\section{Skills}
\skill{Mathematics}{algebraic topology, finite group theory, homotopy theory, representation theory, optimization, numerical analysis, stochastic processes, differential equations, causal inference, statistical modeling.}
\skill{Formal proof and programming}{Lean, Mathlib, Python, PyTorch, NumPy, Pandas, MATLAB, R, C++, LaTeX, Git/GitHub, Linux.}
\skill{Languages}{Mandarin Chinese, English.}

\end{document}
